% !TeX root = ../protokoll.tex
\section{Ziel und Reaktionsübersicht}
\label{sec:ziel}

Ziel dieses Versuchs ist die Herstellung von 2-Methyl-2-butanol durch
säurekatalysierte Hydratisierung von 2-Methyl-2-buten.

\subsection{Reaktionsgleichung}
\begin{reaction}
  \ce{(CH3)2C=CHCH3} + \ce{H2O} ->[\text{H2SO4}] \ce{(CH3)2C(OH)CH2CH3}
  \label{eq:hydrat}
\end{reaction}

\subsection{Stöchiometrie}
\begin{table}[H]
  \centering
  \caption{Eingesetzte Substanzen.}
  \label{tab:stoffe}
  \begin{tabular}{llll}
    \toprule
    Substanz & Menge & $M$ in \si{\gram\per\mol} & Äquivalente \\
    \midrule
    2-Methyl-2-buten & \SI{10.0}{\milli\liter} & 70.13 & 1.00 \\
    \ce{H2SO4} (konz.) & 2 Tropfen & 98.08 & katalytisch \\
    Wasser & \SI{50}{\milli\liter} & 18.02 & Überschuss \\
    \bottomrule
  \end{tabular}
\end{table}

\subsection{Eigenschaften der Zielverbindung}
2-Methyl-2-butanol: $M = \SI{88.15}{\gram\per\mol}$,
$T_{\mathrm{Sdp}} = \SI{102}{\celsius}$,
$\rho = \SI{0.81}{\gram\per\cubic\centi\meter}$, farblose Flüssigkeit.
